\begingroup
\renewcommand{\arraystretch}{1.45}
\fancyhead[L]{}
\fancyhead[C]{\small Das wunderbare Pentagramm}
\par\medskip\Needspace{5\baselineskip}\pdfbookmark[1]{Art. 3}{pentagramma-art-3}\centerline{[3.]}\nobreak\smallskip
Die Beziehungen zwischen den Seiten des sphärischen Fünfecks sind wie folgt
\begingroup\small\setlength{\arraycolsep}{3pt}\[
\begin{array}{cccc} \multicolumn{3}{c}{\text{Quadrate der}}&\text{Gleichungen}\\[9pt] \text{Tangenten}&\text{Sekanten}&\text{Kosekanten}&\\[9pt] \alpha&\gamma\delta&\frac{\gamma\delta}{\alpha}&1+\alpha=\gamma\delta\\[9pt] \beta&\delta\varepsilon&\frac{\delta\varepsilon}{\beta}&1+\beta=\delta\varepsilon\\[9pt] \gamma&\varepsilon\alpha&\frac{\varepsilon\alpha}{\gamma}&1+\gamma=\varepsilon\alpha\\[9pt] \delta&\alpha\beta&\frac{\alpha\beta}{\delta}&1+\delta=\alpha\beta\\[9pt] \varepsilon&\beta\gamma&\frac{\beta\gamma}{\varepsilon}&1+\varepsilon=\beta\gamma \end{array}
\]\endgroup
Diese Gleichungen sind nicht unabhängig; es gilt nämlich identisch:
\[
\begin{gathered} (1+\gamma)(1+\beta-\delta\varepsilon) -(1+\beta)(1+\gamma-\varepsilon\alpha) \\[9pt]\quad{}=\varepsilon\{(1+\beta)\alpha-(1+\gamma)\delta\}\\[9pt] \\[9pt]\quad{}=\varepsilon\{(\alpha\beta-\delta-1)-(\gamma\delta-\alpha-1)\} \end{gathered}
\]
und auf ähnliche Weise wird die fünfte aus dreien der übrigen abgeleitet.

\Needspace{16\baselineskip}
Aus zweien der Größen $\alpha,\beta,\gamma,\delta,\varepsilon$ folgen die übrigen
\par\nopagebreak% portrait-table start
\begingroup\small\setlength{\arraycolsep}{2pt}\[
\begin{array}{llll} \beta=\frac{1+\alpha+\gamma}{\alpha\gamma}, &\beta=\frac{1+\delta}{\alpha}, &\gamma=\frac{1+\alpha}{\alpha\beta-1}, &\beta=\frac{1+\varepsilon}{\alpha\varepsilon-1}\\[9pt] \delta=\frac{1+\alpha}{\gamma}, &\gamma=\frac{1+\alpha}{\delta}, &\delta=\alpha\beta-1, &\gamma=\alpha\varepsilon-1\\[9pt] \varepsilon=\frac{1+\gamma}{\alpha}, &\varepsilon=\frac{1+\alpha+\delta}{\alpha\delta}, &\varepsilon=\frac{1+\beta}{\alpha\beta-1}, &\delta=\frac{1+\alpha}{\alpha\varepsilon-1} \end{array}
\]\endgroup\par% portrait-table end


\Needspace{24\baselineskip}
Beispiel


\SourceLeaf{494}

\[
\begin{array}{lccrrr} &\cos^{2}&\sin^{2}&&&\\[9pt] \alpha=9&\frac{1}{10}&\frac{9}{10}&71^{\circ}&33'&56''\\[9pt] \beta=\frac{3}{2}&\frac{3}{5}&\frac{2}{5}&39&13&54\\[9pt] \gamma=2&\frac{1}{3}&\frac{2}{3}&54&44&7\\[9pt] \delta=5&\frac{1}{6}&\frac{5}{6}&65&54&19\\[9pt] \varepsilon=\frac{1}{3}&\frac{3}{4}&\frac{1}{4}&30&0&0 \end{array}
\]
Schöne Gleichung
\[
\begin{gathered}3+\alpha+\beta+\gamma+\delta+\varepsilon =\alpha\beta\gamma\delta\varepsilon \\[9pt]\quad{}=\surd(1+\alpha)(1+\beta)(1+\gamma)(1+\delta)(1+\varepsilon)\end{gathered}
\]
\par Der Flächeninhalt des sphärischen Pentagons ist $360^{\circ}$ minus die Summe der Seiten.
Setzt man die Summe $=S$ und
\[
(1+i\surd\alpha)(1+i\surd\beta)(1+i\surd\gamma) (1+i\surd\delta)(1+i\surd\varepsilon)=A+Bi
\]
so wird:
\[
\begin{gathered} A=\alpha\beta\gamma\delta\varepsilon.\cos S\\[9pt] B=\alpha\beta\gamma\delta\varepsilon.\sin S \end{gathered}
\]

\Needspace{24\baselineskip}
\par\medskip\Needspace{5\baselineskip}\pdfbookmark[1]{Art. 4}{pentagramma-art-4}\centerline{[4.]}\nobreak\smallskip
\[
\begin{gathered} G(\alpha x+\beta y+\gamma z)^2 +G'(\alpha'x+\beta'y+\gamma'z)^2 \\[9pt]\quad{}+G''(\alpha''x+\beta''y+\gamma''z)^2\\[9pt] \quad=Axx+Byy+Czz+2ayz+2bxz+2cxy \end{gathered}
\]
\[
\begin{array}{rrrr} (A-G)\alpha+&c\beta+&b\gamma&=0\\[9pt] c\alpha+&(B-G)\beta+&a\gamma&=0\\[9pt] b\alpha+&a\beta+&(C-G)\gamma&=0 \end{array}
\]
\[
\begin{gathered}\{a(A-G)-bc\}\alpha\\[9pt]\quad{}=\{b(B-G)-ac\}\beta\\[9pt]\quad{}=\{c(C-G)-ab\}\gamma\end{gathered}
\]
\[
\frac{bc}{a(A-G)-bc}+\frac{ac}{b(B-G)-ac} +\frac{ab}{c(C-G)-ab}+1=0
\]
\[
\begin{gathered}(A-G)(B-G)(C-G)+2abc\\[9pt]\quad{}=aa(A-G)+bb(B-G)+cc(C-G)\end{gathered}
\]
\[
\alpha\alpha=\frac{1}{1+ \left(\frac{a(A-G)-bc}{b(B-G)-ac}\right)^2+ \left(\frac{a(A-G)-bc}{c(C-G)-ab}\right)^2}
\]


\Needspace{22\baselineskip}
\SourceLeaf{495}

Setzt man
\[
\frac{bc}{a(A-x)-bc}+\frac{ac}{b(B-x)-ac}+\frac{ab}{c(C-x)-ab}+1=\Omega
\]
so gilt identisch
\[
\begin{gathered} \{a(A-x)-bc\}\,\{b(B-x)-ac\}\,\{c(C-x)-ab\}\,\Omega\\[9pt] =-abc(x-G)(x-G')(x-G'') \end{gathered}
\]


\Needspace{24\baselineskip}
Differenziert man und setzt nach der Differenziation $x=G$, so wird
\[
\begin{gathered} \{a(A-G)-bc\}\,\{b(B-G)-ac\}\{c(C-G)-ab\}\\[9pt] \times abc\Bigl\{\frac{1}{(a(A-G)-bc)^2}+\frac{1}{(b(B-G)-ac)^2}\\[9pt]\quad{}+\frac{1}{(c(C-G)-ab)^2}\Bigr\}\\[9pt] =abc(G-G')(G-G'') \end{gathered}
\]
woraus
\[
\alpha=\surd{\frac{[b(B-G)-ac][c(C-G)-ab]}{[a(A-G)-bc](G-G')(G-G'')}}
\]

Die Schlüsse müssen abgeändert werden, wenn eine der Größen $a,b,c$
verschwindet. Die obige erste Gleichung für $G$ wäre dann eine Identität.

\clearpage
\endgroup
